La Tabla~\ref{tab:master} recoge, en un único lugar, el estado real de los
24 parámetros originalmente identificados en el reporte anterior (ARC-88),
incluyendo los que este documento y sus extensiones (ARC-89/90/91) no
tocan directamente -- para que ningún parámetro quede implícito o se
confunda ``mencionado en el resumen'' con ``auditado en este documento''.
La columna \textbf{Evidencia} indica dónde vive la justificación real de
cada uno; cuando es este documento, se referencia la sección; cuando es el
reporte anterior (medición de campaña real, no barrido dedicado) o no
existe medición, se indica explícitamente.

\begin{longtable}{@{}p{0.03\textwidth}p{0.24\textwidth}p{0.1\textwidth}p{0.13\textwidth}p{0.44\textwidth}@{}}
\toprule
\textbf{\#} & \textbf{Parámetro} & \textbf{Valor} & \textbf{Estado} & \textbf{Evidencia} \\
\midrule
\endhead
1  & \texttt{interval\_ns} (CPU) & 1\,ms & Medido (este doc.) & Sección~\ref{sec:interval-ns} -- valor razonable y conservador, verificado en esta plataforma \\
2  & \texttt{gpu\_interval\_ns} & 5\,ms & Medido (este doc.) & Sección~\ref{sec:gpu-interval-ns} -- viable y conservador; no aislado como único óptimo frente a 10\,ms \\
3  & \texttt{target\_windows\_per\_repetition} (CPU) & 50 & Medido (este doc.), reencuadrado & Sección~\ref{sec:repetitions} -- guardarraíl de duración mínima, nunca activo; no es prueba de suficiencia estadística \\
4  & \texttt{target\_windows\_per\_repetition} (GPU) & 5 & Medido (este doc.), reencuadrado & Sección~\ref{sec:repetitions} -- misma salvedad que \#3, más la resolución NVML de $\sim$100\,ms (Sección~\ref{sec:gpu-interval-ns}) \\
5  & \texttt{running\_ratio\_min} & 0{,}90 & Medido (reporte anterior) & \emph{Reporte\_Justificacion\_Parametros\_Experimentales.md} §2.3 -- valor real observado = 1{,}0 en el 100\,\% de las corridas; sin barrido dedicado en este documento \\
6  & \texttt{repetitions\_per\_combination} & 10 (final CPU); 3 (smoke) & Medido (este doc.), decisión conservadora & Sección~\ref{sec:repetitions} -- 3 es solo el mínimo operativo; el protocolo final usa el máximo $n$ ensayado porque no hay restricción de cómputo y reduce el riesgo de omitir anomalías tardías \\
7  & \texttt{D03\_TOLERANCE\_FRACTION} (CPU) & 40\,\% & Medido (este doc.), guardarraíl & Sección~\ref{sec:calibration-distribution} -- confirmado conservador frente al ruido real; no derivado estadísticamente \\
8  & Tolerancia cruzada GPU & 5\,\% & Parcial & Sección~\ref{sec:calibration-distribution} -- margen confirmado frente al ruido; el control de reloj GPU fue habilitado después, pero este informe no repite el barrido end-to-end \\
9  & Umbral de estabilidad CAL-10 & 5\,\% CV & Sin justificar de origen & \emph{Reporte\_Justificacion\_Parametros\_Experimentales.md} §2.7 -- razonable en la práctica, sin cita propia; no tocado en este documento \\
10 & \texttt{CV\_THRESHOLD\_PCT} (calentamiento) & 5\,\% & Medido (este doc.) & Sección~\ref{sec:warmup-sensitivity} -- robusto en la mayoría de kernels, con excepciones reales documentadas (\texttt{rodinia\_lud}, \texttt{npb\_cg}); antecedente de Georges et al. recontextualizado (JVM, no HPC nativo) \\
11 & \texttt{MARGIN} (post-detección) & $\times$1{,}2 & Parcial & Sección~\ref{sec:warmup-sensitivity} -- barrido cuantifica el efecto en segundos (0{,}01--0{,}73\,s según kernel), pero sin evidencia de que $\times 1{,}2$ sea suficiente ni excesivo; decisión de ingeniería pendiente de validación estadística entre repeticiones \\
12 & \texttt{min\_mean\_floor} (piso GPU) & 5{,}0\,\% util & Medido (este doc.) & Sección~\ref{sec:warmup-sensitivity} -- confirmado del lado correcto con margen en los 3 kernels GPU elegibles; el codo exacto cae en (2,5]\,\%, no resuelto con mayor precisión \\
13 & \texttt{plateau\_ratio} (changepoint) & 0{,}8 & Parcial, método no activado & Sección~\ref{sec:warmup-sensitivity} -- ninguno de los 6 kernels de la muestra activó el método de respaldo por punto de cambio; el valor sigue respaldado solo por evidencia histórica de ARC-86, barrido dirigido pendiente \\
14 & \texttt{min\_relative\_gain} / \texttt{max\_depth} (changepoint) & 0{,}10 / 6 & \textbf{Sin justificar, NO cubierto por este documento} & \texttt{sensitivity\_warmup\_detector.py} pasa estos valores fijos, nunca los varía -- no confundir con los 4 parámetros sí barridos en \#10-13 \\
15 & \texttt{\_adaptive\_window\_size} & $\max(3,\min(15,n/4))$ & Sin justificar & Heurística propia sin barrido de sensibilidad; no tocado en este documento \\
16 & Niveles de frecuencia F0--F4 & 5 puntos equiespaciados + REF & Pendiente de revalidación CPU sin Turbo & Sección~\ref{sec:literatura} -- la escritura cpufreq funciona, pero el helper administrativo para deshabilitar Turbo aún solicita contraseña en la verificación del 2026-08-14; no se acepta la campaña anterior como evidencia del barrido \\
17 & \texttt{SAFETY\_MARGIN} (timeout) & $\times$3{,}0 & Medido (reporte anterior) & \emph{Reporte\_Justificacion\_Parametros\_Experimentales.md} §2.12 -- presupuesto real usado 1--22\,\%; no tocado en este documento \\
18 & \texttt{timeouts\_seconds} base & ready=15, run=180, shutdown=15 & Sin justificar, bajo impacto & \emph{Reporte\_Justificacion\_Parametros\_Experimentales.md} §2.13 \\
19 & Umbral de carga externa E08 & 1{,}0 & Justificado por convención e integrado & Sección~\ref{sec:literatura} -- el valor no proviene de un barrido, pero el pipeline actual lo aplica antes de calibrar y por combinación, y conserva la observación \\
20 & Rango de temperatura E02 & $[0,90]\,^\circ$C & Parcial e integrado & Sección~\ref{sec:literatura} -- 90\,$^\circ$C no coincide con el $T_{\text{case}}$ de 81\,$^\circ$C del SKU exacto; el sensor ya no se simula: se descubre por \texttt{coretemp}/\texttt{Package id N} y su ausencia bloquea la campaña final CPU \\
21 & \texttt{delegated\_cpus} & 6 núcleos & Parcial & \emph{Reporte\_Justificacion\_Parametros\_Experimentales.md} §2.16 -- topología justificada, el conteo exacto no; no tocado en este documento \\
22 & \texttt{smt\_policy} & 1 hilo/núcleo físico & \textbf{Medido (este doc.)} & Sección~\ref{sec:smt-policy} -- IPC cae 40{,}4\,\% con SMT completo, sin cambio en MPKI/LLC (descarta más tráfico de caché; contención de pipeline es la causa compatible, no la única aislada -- el diseño cambia OMP\_NUM\_THREADS a la vez que SMT) \\
23 & Lanzamientos \texttt{ncu} por kernel & Varía, 5--200 & Medido (este doc.) & Sección~\ref{sec:ncu-convergence} -- criterio de convergencia declarado; \texttt{rodinia\_lud} cerrado con extensión a 200 \\
24 & Peso $w$ del EDP (Fase 4) & $w=1$ primario, $w=2$ sensibilidad & Fijado por diseño (este doc.) & Sección~\ref{sec:literatura} -- pre-registrado antes de Fase 4 para evitar selección posterior de métrica \\
\bottomrule
\caption{Tabla maestra de los 24 parámetros identificados en ARC-88, con su estado real y ubicación de la evidencia tras ARC-89/90/91.}
\label{tab:master}
\end{longtable}

De los 24: \textbf{10 fueron medidos con barrido dedicado en este documento
y sus extensiones} (\#1,2,3,4,6,7,10,12,22,23 -- incluyendo el cierre nuevo
de \texttt{smt\_policy} y de la convergencia de \texttt{rodinia\_lud}; esto
no significa que ninguno tenga matices -- \#3/\#4 quedan reencuadrados como
guardarraíl, no prueba de suficiencia estadística, \#6 documenta un límite
real de detección de anomalías, \#7 es guardarraíl conservador y no
tolerancia derivada, y \#22 tiene causa parcialmente aislada -- ver cada
sección para el matiz exacto), \textbf{2 provienen de medición directa ya
reportada en el documento anterior} sin barrido adicional aquí (\#5,17),
\textbf{5 quedan parcialmente justificados con condición de cierre
declarada} (\#8 tolerancia cruzada GPU, \#11 \texttt{MARGIN} -- cuantificado
pero no validado como suficiente, \#13 \texttt{plateau\_ratio} -- método de
respaldo nunca activado en este barrido, \#20 rango de temperatura E02, y
\#21 \texttt{delegated\_cpus} -- topología justificada, conteo exacto no),
\textbf{2 quedan justificados por decisión metodológica o convención}
(\#16 F0--F4 pendiente de revalidación CPU sin Turbo; \#19 E08 ya integrado
pero sustentado por convención, no por barrido -- ver
Sección~\ref{sec:literatura}), \textbf{1 queda fijado por diseño para
evitar sesgo de selección posterior} (\#24, peso del EDP), y \textbf{4
permanecen explícitamente sin
justificar} (\#9, \#14, \#15, \#18), sin haber sido tocados por este
documento ni por su predecesor -- listados aquí precisamente para que no
desaparezcan de la vista por no formar parte del alcance de ningún barrido
hecho hasta ahora.
